%----------------------------------------------------------------------------------------
%	SECTION TITLE
%----------------------------------------------------------------------------------------

\cvsection{Experience}

%----------------------------------------------------------------------------------------
%	SECTION CONTENT
%	Academia variant: ASSIST Research Lab (the research narrative) kept in
%	full; the two Avidyne industry-engineering roles condensed to a few
%	bullets each so they provide supporting context without crowding out
%	the research/mentoring/funding record. Keep in sync by hand with
%	experience.tex (the industry variant) when facts change -- dates,
%	titles, and metrics must match; only bullet selection differs.
%----------------------------------------------------------------------------------------

\begin{cventries}

%------------------------------------------------
\cventry
{Software Engineer II}
{Software and Systems Engineering, \href{https://www.avidyne.com/}{Avidyne}}
{\textbf{Melbourne, FL}}
{\textbf{May 2026 – Present}}
{
\begin{cvitems}
  \item Design and implement \textbf{embedded C++ software} for avionics flight display systems, spanning \textbf{operating system services}, \textbf{device drivers}, and \textbf{touchscreen-based user interfaces}.
  \item Verify embedded software correctness through \textbf{unit and integration testing}, in accordance with \textbf{FAA} and internal certification standards.
\end{cvitems}
}
\cventry
{Software Engineer Intern / Co-Op}
{Software and Systems Engineering, \href{https://www.avidyne.com/}{Avidyne}}
{\textbf{Melbourne, FL}}
{\textbf{May 2025 – May 2026}}
{
\begin{cvitems}
  \item Expanded \textbf{end-to-end (E2E) test and evaluation (E2TE)} coverage across navigation, communication, flight management, and flight display systems, increasing overall test coverage by \textbf{25\%}.
  \item Developed, integrated, and debugged \textbf{C-based flight software} in the avionics stack, ensuring compliance with real-time, safety-critical \textbf{DO-178C} guidelines for production-ready systems.
  \item Built \textbf{30+ system-level and flight-specific test cases} and \textbf{TFC (Test Framework Core)} software tests, reducing initial verification cycle time by \textbf{10\%}.
\end{cvitems}
}
\cventry
{Research Scholar}
{\href{https://research.fit.edu/assist-lab/}{ASSIST Research Lab}, \href{www.fit.edu}{Florida Institute of Technology}}
{\textbf{Melbourne, FL}}
{\textbf{August 2021 - Present}}
{
\begin{cvitems}
\item Working with a team of research professionals for formal methods of verification and run-time assurance, ML, IoT, robotics, and cyber security.
\item Contributed with \textbf{professors on the development of research proposals} on diverse topics, including cognitive agents on human behavior, the assurance of artificial intelligence in safety-critical systems, and the fine-tuning of LLMs for domain-specific queries.
\item Collaborating on NASA’s \textit{\textbf{University Leadership Initiative (Round 8)}} as part of a \textbf{Florida Tech}-led multi-university/industry team to develop a framework for \textbf{trustworthy, increasingly autonomous aviation safety systems}; partners include \textit{\textbf{Penn State, NC A\&T State, UF, Stanford, Santa Fe College, Uni of New Mexico, Collins Aerospace, and ResilienX}}; part of awards totaling up to \$6.9M over three years.
\item Collaborating with \textit{\textbf{Collins Aerospace, Iowa State University, RTX Technologies Research Center (RTRC), and Smart Information Flow Technologies (SIFT)}}, funded by \textit{\textbf{DARPA}} with the task of formally modeling human cognitive behavior representation with respect to cyber-sickness in AR/VR systems.
\item Collaborated with \textit{\textbf{Penn State University}} on the application of \textbf{Large Language Model (LLM) translation} for cognitive architectures, focusing on enhancing the integration of LLMs to facilitate communication and knowledge transfer within cognitive systems.
\item Contributed with \textit{\textbf{Critical Frequency Design}}, funded by \textit{\textbf{Naval Air Systems Command}}, with the task of developing a modeling approach for designing, maintaining, and supporting air and sea platform fiber optic communications technology.
\item Collaborated with \textit{\textbf{Rockwell Collins and Soar Tech}}, funded by \textit{\textbf{NASA}} with the task of formally verifying the autonomous agent to assure safety as well as the logical correctness of the safety-critical system.
\item Advising a Ph.D. student and undergraduate students on transfer learning, automated data labeling, and assurance frameworks for vision-based classification in autonomous aircraft systems funded by \textit{\textbf{NASA}}, addressing safety and reliability challenges in aviation technologies.
\item Advised \textbf{4 groups of computer science students} on the design, development, and deployment of software for their senior projects
\item Mentoring \textbf{undergraduate} and \textbf{high school students} on machine learning engineering approaches in the aerospace and systems domains, leading to \textbf{conference publications} that addressed real-world challenges in these fields.
\item Investigated the development of a \textbf{cognitive-enhanced agent} for automatically piloting aircraft in dense urban environments which emphasized safe and reliable takeoff/landing among aerial traffic without human intervention.
\item Collected and presented \textbf{AssistTaxi}, a novel dataset for runway and taxiway analysis, contributing to autonomous operations.
\item Developed an \textbf{autonomous aircraft perception system} for accurately detecting and labeling line markings on an airport taxiway.
\item Assisted with the \textbf{formulation of quizzes and homework projects} for the courses: Python, Database Systems, Web Applications, Big Data and Management, and Software Metrics.
\item Recognized with multiple honors and leadership roles, including \textbf{Outstanding Student of the Year} at Honors Convocation 2025, \textbf{Inducted Member of Phi Kappa Phi}, and \textbf{President of the Florida Tech Badminton Club}.
\end{cvitems}
}

\cventry
{Graduate Research Assistant}
{\href{https://research.fit.edu/l3hiai/}{L3Harris Institute for Assured Information}, \href{www.fit.edu}{Florida Institute of Technology}}
{\textbf{Melbourne, FL}}
{\textbf{May 2024 – June 2024}}
{
\begin{cvitems}
\item Developed a \textbf{decentralized framework} enabling \textbf{multiple autonomous agents}—robotic dogs, drones, and mobile robots—to coordinate, communicate, and reach shared goals.
\item Collaborated with \textbf{developers and professors} to rigorously test the system in both \textbf{simulation} and \textbf{real-world environments}.
\end{cvitems}
}

\cventry
{IRI Research Developer}
{\href{https://www.fit.edu/faculty-profiles/s/slhoub-khaled-/}{IRI Research},  \href{www.fit.edu}{Florida Institute of Technology}}
{\textbf{Melbourne, FL}}
{\textbf{May 2023 - July 2023}}
{
\begin{cvitems}
\item Proposed and \textbf{implemented a framework using AI language models} to automatically extract software requirements from source code.
\item Supervised and coordinated with undergraduate students towards the development process.
\end{cvitems}
}
%------------------------------------------------

\end{cventries}
